\section{One residue pair through six models}

Return to the positions $(i,j)$ used throughout the book.  Instead of asking which model is most
powerful, ask what object each model assigns to this pair and what operation would be required to
compare the assignments.

\subsection*{Portrait I: the PSSM has no edge}

A position-specific scoring matrix assigns a field $h_i(a)$ to each residue identity at each site.
Its reconstruction rule may be written
\begin{equation}
P(X_i=a)=\operatorname{softmax}_a h_i(a).
\end{equation}
There is no off-diagonal operator from $j$ to $i$.  Correlation can still appear in the observed
family because sequences are phylogenetically related or sampled from subfamilies, but the PSSM
does not represent that correlation as a pair interaction.  In operator coordinates it occupies
the node-field branch and sets
\begin{equation}
\mathcal K_{i\leftarrow j}=0.
\end{equation}
This apparently trivial portrait is useful because it prevents every successful residue predictor
from being described as an interaction model.

\subsection*{Portrait II: the Potts model assigns a static categorical operator}

A pairwise Potts model attaches a table $J_{ij}(a,b)$ to the unordered pair.  After reference
projection,
\begin{equation}
G_{ij}^{\mathrm P}=P_iJ_{ij}P_j^{\mathsf T}
\end{equation}
is a static operator shared by every sequence evaluated under that fitted family model.  The
conditional probability at $i$ changes with the surrounding sequence because all neighboring
terms are summed, but the response rule contributed by $j$ is fixed.

The edge belongs to one joint Gibbs law when the pair tables and fields are assembled consistently.
That gives the Potts block a stronger probability semantics than an attention route.  It does not
give it physical units, remove phylogenetic bias, or make every nonzero block a contact.  Those are
estimation and validation questions.

\subsection*{Portrait III: Gaussian precision changes the state space}

For a nondegenerate Gaussian vector, a zero precision block has an exact conditional-independence
meaning.  The edge object is a linear map between continuous coordinates, not a table between
categorical contrasts.  If one-hot residue variables are approximated by a Gaussian, the precision
matrix may provide a useful estimator, but the approximation does not erase the simplex constraint
or reference-centered categorical coordinates.

The correspondence is therefore structural:
\begin{equation}
\text{precision block}
\longleftrightarrow
\text{conditional linear edge},
\end{equation}
not literal equality with a Potts coupling.  The two models agree that conditional dependence is
encoded by an off-diagonal block; they disagree about the node space and probability law in which
that block lives.

\subsection*{Portrait IV: Hopfield--Potts compresses edges into collective patterns}

A Hopfield--Potts model factorizes static couplings through patterns
\begin{equation}
J_{ij}(a,b)=\sum_{\mu=1}^{p}\lambda_\mu
\xi_{i\mu}(a)\xi_{j\mu}(b).
\end{equation}
For fixed $p$, many edge operators share a small collection of categorical directions.  The model
is dense in position support but low rank in its global parameterization.  One pattern can touch
many positions without declaring every touched pair to be a separate physical contact.

This portrait separates support from rank.  A sparse Potts graph may carry independent full-rank
operators on a few edges.  A Hopfield--Potts model may carry correlated low-rank operators on nearly
all edges.  Both are pairwise energies; their compression assumptions are different.

\subsection*{Portrait V: an RBM stores motifs and generates all visible orders}

A restricted Boltzmann machine introduces hidden variables $r_\mu$ coupled to visible residues.
Conditional on the hidden layer, the visible sites separate.  After the hidden variables are
marginalized, each softplus term depends jointly on every visible position connected to that hidden
unit.  Its expansion contains one-body, pairwise, and higher-order components.

An RBM hidden unit is therefore closer to a distributed motif than to one edge.  Its pairwise
operator can be extracted by functional projection, but that operator is only one shadow of the
latent factor.  Reading all pairs touched by a hidden unit as direct contacts would confuse latent
support with visible interaction order.

The contrast with modern Hopfield retrieval is equally important.  Independent Bernoulli hidden
variables allow several motifs to be active simultaneously.  A categorical Hopfield softmax makes
stored memories compete on one simplex.  Both use log-sum-exp after marginalization; they sum over
different latent sample spaces.

\subsection*{Portrait VI: the Transformer has a route and an audited response}

For a Transformer head, the native edge coefficient is
\begin{equation}
\alpha_{ij}^{(h)}(x)
=\operatorname{softmax}_j s_{ij}^{(h)}(x).
\end{equation}
It tells us how a target distributes routing mass over source positions.  The associated frozen
value-path map lies in the span of the head value/output maps.  Both quantities depend on the
current hidden state and belong to the computation.

The operator comparable with a Potts block is obtained elsewhere: vary the source residue, hold a
declared background fixed, observe the categorical output at the target, and project the resulting
table.  This gives $G_{i\leftarrow j}(C)$, a directed and generally background-dependent response.
The route and response may correlate, but they answer different questions:
\begin{equation}
\begin{aligned}
\alpha_{ij}^{(h)}(x)&:\quad \text{where one layer reads},\\
G_{i\leftarrow j}(C)&:\quad \text{how the complete model prediction changes}.
\end{aligned}
\end{equation}

The Transformer portrait is consequently double.  It is a dynamic graph computation internally
and a background-dependent operator field under external audit.  Neither object is automatically a
joint sequence energy.

\section{What survives comparison}

The six portraits meet only after a typed extraction.  PSSM fields, Potts blocks, Gaussian
precision, Hopfield factors, RBM projections, attention routes, and pLM responses should not be
placed in one matrix merely because they share position indices.  What can be compared is more
specific:
\begin{equation}
\boxed{
\text{state space}
+\text{edge operator}
+\text{context}
+\text{order}
+\text{normalization}
+\text{estimator}.}
\end{equation}
Agreement in the edge-operator coordinate can then be measured while the remaining differences are
kept as part of the result.
